Localizing objects and parts from natural language in 3D space is essential for robotics, AR, and embodied AI, yet existing methods face a trade-off between the accuracy and geometric consistency of per-scene optimization and the efficiency of feed-forward inference. We present TrianguLang, a feed-forward framework for 3D localization that requires no camera calibration at inference. Unlike prior methods that treat views independently, we introduce Geometry-Aware Semantic Attention (GASA), which utilizes predicted geometry to gate cross-view feature correspondence, suppressing semantically plausible but geometrically inconsistent matches without requiring ground-truth poses. Validated on five benchmarks including ScanNet++ and uCO3D, TrianguLang achieves state-of-the-art feed-forward text-guided segmentation and localization, reducing user effort from $O(N)$ clicks to a single text query. The model processes each frame at 1008x1008 resolution in $\sim$57ms ($\sim$18 FPS) without optimization, enabling practical deployment for interactive robotics and AR applications. Code and checkpoints are available at \url{https://cwru-aism.github.io/triangulang/}.


% We present TrianguLang, a feed-forward framework for part-level 3D localization that requires no camera calibration at inference. Our key insight is that predicted geometry from foundation models can serve as an "arbiter" for semantic attention, suppressing false correspondences that arise when views disagree. We introduce Geometry-Aware Semantic Attention (GASA), a novel attention mechanism that combines semantic similarity with geometric proximity, attending only to features that are both semantically similar and geometrically close in 3D space. Combined with world-space positional encoding and spatial language parsing, GASA enables zero-shot transfer to novel scenes without camera calibration. Across five benchmarks: ScanNet++, uCO3D, NVOS, SPIn-NeRF, and PartImageNet; Triangulang achieves state-of-the-art performance on text-guided multi-view segmentation while reducing user effort from $O(N)$ clicks to a single text query. TrianguLang operates at $\sim$10 FPS with no per-scene optimization, enabling practical deployment for interactive robotics and AR applications.